% !TEX program = xelatex
\documentclass[UTF8,fontset=fandol,a4paper,11pt]{ctexart}
\usepackage[margin=21mm,top=23mm,bottom=22mm,headheight=16pt]{geometry}
\usepackage{amsmath,amssymb,mathtools,bm}
\usepackage{tikz}
\usetikzlibrary{arrows.meta,calc,angles,quotes,patterns}
\usepackage{xcolor,booktabs,array,longtable,enumitem,fancyhdr,needspace}
\usepackage[unicode,hidelinks]{hyperref}
\definecolor{ink}{HTML}{24343C}
\definecolor{accent}{HTML}{166A70}
\definecolor{muted}{HTML}{65747A}
\definecolor{rulegray}{HTML}{D4DDE0}
\definecolor{warm}{HTML}{966037}
\color{ink}
\setlength{\parindent}{2em}
\setlength{\parskip}{4pt}
\linespread{1.12}
\AtBeginDocument{\setlength{\abovedisplayskip}{6pt plus 2pt minus 1pt}
\setlength{\belowdisplayskip}{6pt plus 2pt minus 1pt}
\setlength{\abovedisplayshortskip}{0pt plus 1pt}
\setlength{\belowdisplayshortskip}{4pt plus 1pt}}
\setlist[enumerate]{label=(\arabic*),leftmargin=2.3em,itemsep=3pt,topsep=4pt}
\setlist[itemize]{leftmargin=2em,itemsep=2pt,topsep=4pt}
\allowdisplaybreaks[1]
\setlength{\emergencystretch}{2em}
\setcounter{secnumdepth}{0}
\providecommand{\tightlist}{\setlength{\itemsep}{0pt}\setlength{\parskip}{0pt}}
\ctexset{section={format=\Large\bfseries\color{accent},beforeskip=10pt,afterskip=9pt},
subsection={format=\normalsize\bfseries\color{accent},beforeskip=10pt,afterskip=4pt},
subsubsection={format=\normalsize\bfseries,beforeskip=8pt,afterskip=3pt}}
\pagestyle{fancy}
\fancyhf{}
\fancyhead[L]{\small\color{muted}物理学难题集萃}
% 页眉仅保留书名；版本名称只在封面显示。
\fancyfoot[L]{\small\color{muted}\chapterlabel}
\fancyfoot[R]{\small\thepage}
\renewcommand{\headrulewidth}{0.4pt}
\renewcommand{\footrulewidth}{0pt}
\newcommand{\dd}{\mathop{}\!\mathrm{d}}
\newcommand{\vect}[1]{\bm{#1}}
\newcommand{\tagline}[1]{\noindent{\small\color{muted}#1}\par}
\newcommand{\result}[1]{\boxed{#1}}
\newcommand{\figcaption}[1]{\par\vspace{2pt}{\small\color{muted}#1}\par}
\newcommand{\teach}[2]{\Needspace{4\baselineskip}\subsection*{讲解要点与追问}#1\par\smallskip\noindent\textbf{追问：}#2\par}
\newcommand{\workarea}{\par\vspace{7mm}\noindent{\small\color{muted}作答区}\par\vspace{4mm}
\noindent\textcolor{rulegray}{\rule{\linewidth}{0.3pt}}\par\vspace{16mm}
\noindent\textcolor{rulegray}{\rule{\linewidth}{0.3pt}}\par\vspace{16mm}
\noindent\textcolor{rulegray}{\rule{\linewidth}{0.3pt}}\par\vspace{16mm}
\noindent\textcolor{rulegray}{\rule{\linewidth}{0.3pt}}}
\tikzset{axis/.style={-{Stealth[length=2mm]},thin,color=muted},
motion/.style={-{Stealth[length=2.3mm]},thick,color=accent},
force/.style={-{Stealth[length=2.3mm]},thick,color=warm},
traj/.style={thick,color=accent},guide/.style={dashed,thin,color=muted},
object/.style={thick,color=ink},every node/.style={font=\small}}\widowpenalty=10000\clubpenalty=10000\displaywidowpenalty=10000

\newcommand{\chapterlabel}{热学第四章\quad 实际气体、液体、固体与相变}
\hypersetup{pdftitle={热学第四章 实际气体、液体、固体与相变·学生版},pdfauthor={},pdfsubject={热学第四章第1--17题}}
\begin{document}
\pdfbookmark[0]{热学第四章 实际气体、液体、固体与相变}{cover}
\thispagestyle{empty}\vspace*{20mm}
\noindent{\Huge\bfseries 物理学难题集萃}\par\vspace{8mm}
\noindent{\LARGE 热学第四章}\par\vspace{5mm}
\noindent{\LARGE 实际气体、液体、固体与相变}\par\vspace{10mm}
\noindent{\Large\color{accent}学生版}\par\vspace{8mm}
\noindent\textcolor{accent}{\rule{\linewidth}{1pt}}\par\vspace{8mm}
\noindent 第1--17题\par\vspace{5mm}
\noindent 题目、题图与作答区。\par\vfill\clearpage
\pdfbookmark[1]{题目目录}{contents}\section*{题目目录}
\begin{longtable}{@{}>{\color{accent}\bfseries}r p{120mm}@{}}
\toprule 题号 & 题目\\\midrule
1 & 范德瓦耳斯气体自由膨胀的熵增 \\
2 & 范德瓦耳斯氧气节流过程温度变化 \\
3 & 乙醚临界现象中气液质量比 \\
4 & 毛细现象永动机方案的热力学破绽 \\
5 & 毛细管中不同长度水柱的液面形状 \\
6 & 双液体U形毛细管表面张力比 \\
7 & 微小水滴的曲率压强与密度变化 \\
8 & 相通肥皂泡的稳定与不稳定平衡 \\
9 & 润湿浮球温度变化时浸入体积不变条件 \\
10 & 对称连体液膜中环形细线张力 \\
11 & 润湿竖直平板上的毛细上升高度 \\
12 & 一维晶体作用势、弹性和断裂形变 \\
13 & NaCl晶体结合能的估算 \\
14 & 等温压缩氟利昂的液化质量 \\
15 & 互不相溶液体共沸温度与蒸汽组成 \\
16 & 冰水水蒸气三相定容加热后的质量分配 \\
17 & 地幔硅酸盐熔点随深度变化 \\\bottomrule\end{longtable}

\clearpage

\pdfbookmark[2]{第 1 题 范德瓦耳斯气体自由膨胀的熵增}{q1}

\section*{第 1 题\quad 范德瓦耳斯气体自由膨胀的熵增}

1 mol范德瓦耳斯气体初温$T_1$、初体积$V_1$，在绝热刚性容器内向真空自由膨胀到体积$V_2>V_1$。已知摩尔定体热容量$C_V$及参数$a,b$均为常量，状态方程为$p=RT/(V-b)-a/V^2$。求终温及熵增，并说明实际不可逆过程和计算熵变所用可逆路径的区别。设初、终态均为稳定单相气体，所用温度、体积落在模型适用范围。

\begin{center}\begin{minipage}{0.96\linewidth}\centering\begin{tikzpicture}[x=1cm,y=1cm,font=\small,every node/.append style={align=center}]\draw[object] (0,0) rectangle (6,2);\fill[blue!12] (0,0) rectangle (2.5,2);
\draw[object,dashed] (2.5,0)--(2.5,2);\node at (1.25,1){气体\\$T_1,V_1$};\node at (4.2,1.4){初始真空};
\draw[motion] (2.8,.6)--(5,.6);\node[below] at (3,-.1){刚性绝热外壁；移去隔板，终体积$V_2$};\end{tikzpicture}\figcaption{图 热4.1　向真空自由膨胀，外界没有活塞功。}\end{minipage}\end{center}

\workarea

\clearpage

\pdfbookmark[2]{第 2 题 范德瓦耳斯氧气节流过程温度变化}{q2}

\section*{第 2 题\quad 范德瓦耳斯氧气节流过程温度变化}

1 mol氧气经绝热多孔塞节流，摩尔体积由高压侧$V_1=4.0\times10^{-3}\,\mathrm{m^3/mol}$增至低压侧$V_2=1.2\times10^{-2}\,\mathrm{m^3/mol}$。范德瓦耳斯参数$a=0.138\,\mathrm{Pa\cdot m^6/mol^2}$、$b=3.2\times10^{-5}\,\mathrm{m^3/mol}$，$C_V=20.8\,\mathrm{J/(mol\cdot K)}$恒定。流动动能和位能变化忽略。求温度变化：先写保留$b$、以初温$T_1$表示的关系，再在可忽略排斥体积对温差影响的近似下计算数值，并判断该近似需要什么条件。

\begin{center}\begin{minipage}{0.96\linewidth}\centering\begin{tikzpicture}[x=1cm,y=1cm,font=\small,every node/.append style={align=center}]\draw[object] (0,0) rectangle (7,1.8);\fill[pattern=north east lines] (3,0) rectangle (3.45,1.8);
\node at (1.4,1.2){高压侧\\$p_1,T_1,V_1$};\node at (5.4,1.2){低压侧\\$p_2,T_2,V_2$};
\draw[motion] (.4,.45)--(2.6,.45);\draw[motion] (3.8,.45)--(6.5,.45);\node[below] at (3.2,-.15){绝热多孔塞};\end{tikzpicture}\figcaption{图 热4.2　气体持续通过多孔塞，两侧流动带有压力功。}\end{minipage}\end{center}

\workarea

\clearpage

\pdfbookmark[2]{第 3 题 乙醚临界现象中气液质量比}{q3}

\section*{第 3 题\quad 乙醚临界现象中气液质量比}

纯乙醚封装在玻璃管内，管容积恰等于所装乙醚的总临界体积。加热时气、液两相共存，达到临界点时界面消失。其摩尔状态方程为$(p+a/V^2)(V-b)=RT$。已知摩尔临界体积$V_K$；温度$T$时，饱和气相与液相的摩尔体积分别为$V_g,V_l$。求两相各占总体积的比例$B_g,B_l$及两相质量比例。在$V_g\gg V_l$时，用$a,b,T$近似表示$V_l$。若20℃封装，取$T_K=466.95\,\mathrm K$、$p_K=34.60\,\mathrm{atm}$，求所需液相体积比例，并核查近似的精度。

\begin{center}\begin{minipage}{0.96\linewidth}\centering\begin{tikzpicture}[x=1cm,y=1cm,font=\small,every node/.append style={align=center}]\draw[object] (0,0) rectangle (1.7,3.4);\fill[blue!20] (0,0) rectangle (1.7,1.05);
\draw[object] (0,1.05)--(1.7,1.05);\node at (.85,.5){液相};\node at (.85,2.2){气相};
\node[right] at (2,2.9){总量$\nu$\\管容积$\nu V_K$};\draw[motion] (3,.5)--(3,1.7) node[above]{缓慢加热};\end{tikzpicture}\figcaption{图 热4.3　两相共存时摩尔数分配与体积分配不同。}\end{minipage}\end{center}

\workarea

\clearpage

\pdfbookmark[2]{第 4 题 毛细现象永动机方案的热力学破绽}{q4}

\section*{第 4 题\quad 毛细现象永动机方案的热力学破绽}

图示密闭装置与温度恒定的大气导热。A为完全润湿的毛细区域，B为实心隔挡，其外侧空间C较大；导管D内有可无摩擦移动的活塞E。活门F在一定正压差下打开、压差近零时关闭；G的开闭顺序相反。所有空白空间只含该液体蒸汽，底部及毛细区装同一液体。设计者假定毛细凹液面附近的饱和蒸汽压$p_1$与下方平液面附近的$p_2$相同，又用蒸汽的重力压差，声称能循环地从单一热源取热并全部变为活塞功。先说明这一推证的过程，再否定它并导出$p_1,p_2$的定量关系。设毛细半径$r$，升高$h$，液体与蒸汽密度为$\rho_l,\rho_v$，表面张力系数为$\sigma$；采用密度近似恒定的小高度模型。

\begin{center}\begin{minipage}{.96\linewidth}\centering
\begin{tikzpicture}[x=.85cm,y=.85cm,font=\small]
\draw[object] (0,0) rectangle (8,6);
\fill[blue!12] (0,0) rectangle (8,1.2);
\fill[blue!12] (0,1.2)--(1.8,1.2)--(1.8,3.2) .. controls (1.2,2.94) and (.6,2.94) .. (0,3.2)--cycle;
\fill[pattern=north east lines] (1.8,1.2) rectangle (4.8,3.2);
\draw[object] (1.8,1.2) rectangle (4.8,3.2);
\node[fill=white,inner sep=2pt] at (3.3,2.1){B：实心隔挡};\node at (.85,1.9){A};
\draw[traj] (0,3.2) .. controls (.6,2.94) and (1.2,2.94) .. (1.8,3.2);
\draw[object] (1.8,4)--(1.8,6) (1.8,4)--(3.6,4) (4.2,4)--(4.8,4);
\draw[object,very thick] (4.8,3.2)--(4.8,4);
\draw[object] (4.8,3.6)--(8.7,3.6);
\draw[object] (1.8,4)--(2.15,3.57);\node[left] at (1.75,3.75){F};
\draw[object] (4.2,4)--(3.85,3.62);\node[above] at (3.9,4.08){G};
\node at (2.95,3.6){D};\node[above right] at (4.8,4.03){E};
\node at (5.7,5.0){C：蒸汽空间};
\fill (.7,3.3) circle (1.3pt) node[above]{$p_1$};
\fill (5.8,3.3) circle (1.3pt) node[above]{$p_0$};
\fill (5.8,1.3) circle (1.3pt) node[above]{$p_2$};
\draw[guide] (5.4,3.2)--(7.5,3.2) (5.4,1.2)--(7.5,1.2);
\draw[<->] (7.2,1.2)--(7.2,3.2) node[midway,right]{$h$};
\draw[object] (0,0) rectangle (8,6);
\end{tikzpicture}\figcaption{图 热4.4　F连通A与活塞左腔D，G连通D与C；活塞右侧属于C。阀片位置仅示意通路。}\end{minipage}\end{center}

\workarea

\clearpage

\pdfbookmark[2]{第 5 题 毛细管中不同长度水柱的液面形状}{q5}

\section*{第 5 题\quad 毛细管中不同长度水柱的液面形状}

两端开口的竖直圆柱毛细管内径为1.0 mm，完全润湿水。管内水柱的高度分别为2.00 cm、4.00 cm和2.98 cm，求各情形下下液面是向液体内凹、向外凸，还是近似平面。取表面张力$\sigma=0.073\,\mathrm{N/m}$、水密度$\rho=10^3\,\mathrm{kg/m^3}$、$g=9.80\,\mathrm{m/s^2}$。上液面按半球处理；下液面允许接触线钉扎或端口约束，以有符号曲率半径描述。

\begin{center}\begin{minipage}{0.96\linewidth}\centering\begin{tikzpicture}[x=1cm,y=1cm,font=\small,every node/.append style={align=center}]\foreach \x/\h/\lab in {0/1.8/2.00,2.6/2.9/4.00,5.2/2.35/2.98}{
\draw[object] (\x,0)--(\x,3.3);\draw[object] (\x+.8,0)--(\x+.8,3.3);
\draw[traj] (\x,3) .. controls (\x+.25,2.7) and (\x+.55,2.7) .. (\x+.8,3);
\draw[guide] (\x,3-\h)--(\x+.8,3-\h);\node[below] at (\x+.4,-.2){$h=\lab\,\mathrm{cm}$};}
\node at (3,4){上液面为凹半球；虚线处下液面形状待求};\end{tikzpicture}\figcaption{图 热4.5　在同一符号约定下比较三个高度。}\end{minipage}\end{center}

\workarea

\clearpage

\pdfbookmark[2]{第 6 题 双液体U形毛细管表面张力比}{q6}

\section*{第 6 题\quad 双液体U形毛细管表面张力比}

总长$L$、内半径$R$均匀的倒U形毛细管，其两开口竖直向下，略浸入两种液体的等高液面。两液体密度$\rho_1,\rho_2$，与管壁接触角均为零；毛细管中气体不漏出，压缩过程等温，液体上升高度分别为$h_1,h_2$。插入前气体压强为大气压$p_0$，开口浸入深度忽略。求两液体表面张力系数$\sigma_1/\sigma_2$，并说明$h_1+h_2<L$的意义。

\begin{center}\begin{minipage}{0.96\linewidth}\centering\begin{tikzpicture}[x=1cm,y=1cm,font=\small,every node/.append style={align=center}]\fill[blue!15] (-.5,-.7) rectangle (1.8,.8);\fill[brown!15] (3.2,-.7) rectangle (5.5,.8);
\draw[object] (-.5,.8)--(1.8,.8);\draw[object] (3.2,.8)--(5.5,.8);
\draw[object] (.5,0)--(.5,3) .. controls (.5,3.5) and (4.5,3.5) ..(4.5,3)--(4.5,0);
\draw[object] (.8,0)--(.8,2.8) ..controls(.8,3.15) and(4.2,3.15)..(4.2,2.8)--(4.2,0);
\draw[traj] (.5,2.1)..controls(.6,1.95) and(.7,1.95)..(.8,2.1);
\draw[traj] (4.2,1.8)..controls(4.3,1.65) and(4.4,1.65)..(4.5,1.8);
\draw[<->] (1.3,.8)--(1.3,2.1) node[midway,right]{$h_1$};\draw[<->] (3.7,.8)--(3.7,1.8) node[midway,left]{$h_2$};
\node at (2.5,3.8){封闭气体$p$，长度$L-h_1-h_2$};\node at (.3,-.45){液体1};\node at (4.6,-.45){液体2};\end{tikzpicture}\figcaption{图 热4.6　倒U形管中的气体被两液柱共同压缩。}\end{minipage}\end{center}

\workarea

\clearpage

\pdfbookmark[2]{第 7 题 微小水滴的曲率压强与密度变化}{q7}

\section*{第 7 题\quad 微小水滴的曲率压强与密度变化}

4℃、1 atm时水的表面张力$\sigma=7.24\times10^{-2}\,\mathrm{N/m}$，等温压缩系数$\beta=-(1/V)(\partial V/\partial p)_T=4.75\times10^{-5}\,\mathrm{atm^{-1}}$。同温、同外压下，平液面附近水的密度为$\rho_0$。求半径$r=1.0\times10^{-6}\,\mathrm{cm}$的微小水滴密度。计算中取$\sigma,\beta$为常量，并讨论该连续介质估算的条件。

\begin{center}\begin{minipage}{0.96\linewidth}\centering\begin{tikzpicture}[x=1cm,y=1cm,font=\small,every node/.append style={align=center}]\fill[blue!12] (0,0) circle (1.2);\draw[object] (0,0) circle (1.2);\draw[guide] (0,0)--(1.2,0) node[midway,above]{$r$};
\node at (0,-.45){$p,\rho$};\node at (2.2,1.2){外压$p_0$};\draw[force] (1.2,0)--(.84,0);
\draw[force] (-1.2,0)--(-.84,0);\node at (3,-.5){$p-p_0=2\sigma/r$};
\node[below] at (0,-1.22){箭头为张力在小曲面上的合力};\end{tikzpicture}\figcaption{图 热4.7　球形水滴只有一个界面，外压不是滴内压强。}\end{minipage}\end{center}

\workarea

\clearpage

\pdfbookmark[2]{第 8 题 相通肥皂泡的稳定与不稳定平衡}{q8}

\section*{第 8 题\quad 相通肥皂泡的稳定与不稳定平衡}

两相同球冠形肥皂泡A、B固定在半径$r_1$相同的圆管口上；外界压强相同，两泡连通，外部供气阀K关闭，连接泡的阀门保持开启。表面张力系数恒定，忽略液膜厚度和气体密度变化。证明：两泡均小于半球时等形平衡稳定，均大于半球时等形平衡不稳定。后一情形取$r_1=2.00\,\mathrm{cm}$、初始球面曲率半径$r_2=2.50\,\mathrm{cm}$，求一泡变小、另一泡变大后再次平衡的曲率半径与球冠高度。

\begin{center}\begin{minipage}{0.96\linewidth}\centering\begin{tikzpicture}[x=1cm,y=1cm,font=\small,every node/.append style={align=center}]\draw[traj] (-0.7500,0.0000) -- (-0.7389,0.0201) -- (-0.7273,0.0399) -- (-0.7151,0.0594) -- (-0.7025,0.0786) -- (-0.6893,0.0974) -- (-0.6756,0.1158) -- (-0.6614,0.1339) -- (-0.6467,0.1516) -- (-0.6316,0.1688) -- (-0.6160,0.1857) -- (-0.6000,0.2021) -- (-0.5835,0.2181) -- (-0.5666,0.2336) -- (-0.5492,0.2487) -- (-0.5315,0.2633) -- (-0.5134,0.2774) -- (-0.4949,0.2911) -- (-0.4761,0.3042) -- (-0.4569,0.3168) -- (-0.4373,0.3289) -- (-0.4175,0.3404) -- (-0.3973,0.3514) -- (-0.3769,0.3619) -- (-0.3561,0.3718) -- (-0.3352,0.3811) -- (-0.3139,0.3899) -- (-0.2925,0.3981) -- (-0.2708,0.4057) -- (-0.2489,0.4127) -- (-0.2269,0.4192) -- (-0.2047,0.4250) -- (-0.1823,0.4302) -- (-0.1598,0.4348) -- (-0.1372,0.4389) -- (-0.1145,0.4423) -- (-0.0917,0.4450) -- (-0.0688,0.4472) -- (-0.0459,0.4488) -- (-0.0230,0.4497) -- (0.0000,0.4500) -- (0.0230,0.4497) -- (0.0459,0.4488) -- (0.0688,0.4472) -- (0.0917,0.4450) -- (0.1145,0.4423) -- (0.1372,0.4389) -- (0.1598,0.4348) -- (0.1823,0.4302) -- (0.2047,0.4250) -- (0.2269,0.4192) -- (0.2489,0.4127) -- (0.2708,0.4057) -- (0.2925,0.3981) -- (0.3139,0.3899) -- (0.3352,0.3811) -- (0.3561,0.3718) -- (0.3769,0.3619) -- (0.3973,0.3514) -- (0.4175,0.3404) -- (0.4373,0.3289) -- (0.4569,0.3168) -- (0.4761,0.3042) -- (0.4949,0.2911) -- (0.5134,0.2774) -- (0.5315,0.2633) -- (0.5492,0.2487) -- (0.5666,0.2336) -- (0.5835,0.2181) -- (0.6000,0.2021) -- (0.6160,0.1857) -- (0.6316,0.1688) -- (0.6467,0.1516) -- (0.6614,0.1339) -- (0.6756,0.1158) -- (0.6893,0.0974) -- (0.7025,0.0786) -- (0.7151,0.0594) -- (0.7273,0.0399) -- (0.7389,0.0201) -- (0.7500,0.0000);\draw[object] (-0.75,0)--(-.2,0) (.2,0)--(.75,0);\draw[traj] (2.2500,0.0000) -- (2.2611,0.0201) -- (2.2727,0.0399) -- (2.2849,0.0594) -- (2.2975,0.0786) -- (2.3107,0.0974) -- (2.3244,0.1158) -- (2.3386,0.1339) -- (2.3533,0.1516) -- (2.3684,0.1688) -- (2.3840,0.1857) -- (2.4000,0.2021) -- (2.4165,0.2181) -- (2.4334,0.2336) -- (2.4508,0.2487) -- (2.4685,0.2633) -- (2.4866,0.2774) -- (2.5051,0.2911) -- (2.5239,0.3042) -- (2.5431,0.3168) -- (2.5627,0.3289) -- (2.5825,0.3404) -- (2.6027,0.3514) -- (2.6231,0.3619) -- (2.6439,0.3718) -- (2.6648,0.3811) -- (2.6861,0.3899) -- (2.7075,0.3981) -- (2.7292,0.4057) -- (2.7511,0.4127) -- (2.7731,0.4192) -- (2.7953,0.4250) -- (2.8177,0.4302) -- (2.8402,0.4348) -- (2.8628,0.4389) -- (2.8855,0.4423) -- (2.9083,0.4450) -- (2.9312,0.4472) -- (2.9541,0.4488) -- (2.9770,0.4497) -- (3.0000,0.4500) -- (3.0230,0.4497) -- (3.0459,0.4488) -- (3.0688,0.4472) -- (3.0917,0.4450) -- (3.1145,0.4423) -- (3.1372,0.4389) -- (3.1598,0.4348) -- (3.1823,0.4302) -- (3.2047,0.4250) -- (3.2269,0.4192) -- (3.2489,0.4127) -- (3.2708,0.4057) -- (3.2925,0.3981) -- (3.3139,0.3899) -- (3.3352,0.3811) -- (3.3561,0.3718) -- (3.3769,0.3619) -- (3.3973,0.3514) -- (3.4175,0.3404) -- (3.4373,0.3289) -- (3.4569,0.3168) -- (3.4761,0.3042) -- (3.4949,0.2911) -- (3.5134,0.2774) -- (3.5315,0.2633) -- (3.5492,0.2487) -- (3.5666,0.2336) -- (3.5835,0.2181) -- (3.6000,0.2021) -- (3.6160,0.1857) -- (3.6316,0.1688) -- (3.6467,0.1516) -- (3.6614,0.1339) -- (3.6756,0.1158) -- (3.6893,0.0974) -- (3.7025,0.0786) -- (3.7151,0.0594) -- (3.7273,0.0399) -- (3.7389,0.0201) -- (3.7500,0.0000);\draw[object] (2.25,0)--(2.8,0) (3.2,0)--(3.75,0);\draw[traj] (-0.7500,-3.0000) -- (-0.7800,-2.9576) -- (-0.8076,-2.9137) -- (-0.8327,-2.8683) -- (-0.8552,-2.8215) -- (-0.8752,-2.7736) -- (-0.8924,-2.7247) -- (-0.9069,-2.6749) -- (-0.9187,-2.6243) -- (-0.9276,-2.5732) -- (-0.9337,-2.5217) -- (-0.9369,-2.4699) -- (-0.9373,-2.4180) -- (-0.9348,-2.3662) -- (-0.9294,-2.3146) -- (-0.9212,-2.2633) -- (-0.9101,-2.2126) -- (-0.8963,-2.1626) -- (-0.8797,-2.1135) -- (-0.8604,-2.0653) -- (-0.8385,-2.0182) -- (-0.8140,-1.9725) -- (-0.7871,-1.9282) -- (-0.7577,-1.8854) -- (-0.7260,-1.8443) -- (-0.6920,-1.8051) -- (-0.6560,-1.7677) -- (-0.6179,-1.7325) -- (-0.5780,-1.6994) -- (-0.5362,-1.6685) -- (-0.4929,-1.6400) -- (-0.4480,-1.6140) -- (-0.4017,-1.5904) -- (-0.3543,-1.5695) -- (-0.3057,-1.5512) -- (-0.2562,-1.5357) -- (-0.2059,-1.5229) -- (-0.1550,-1.5129) -- (-0.1036,-1.5057) -- (-0.0519,-1.5014) -- (0.0000,-1.5000) -- (0.0519,-1.5014) -- (0.1036,-1.5057) -- (0.1550,-1.5129) -- (0.2059,-1.5229) -- (0.2562,-1.5357) -- (0.3057,-1.5512) -- (0.3543,-1.5695) -- (0.4017,-1.5904) -- (0.4480,-1.6140) -- (0.4929,-1.6400) -- (0.5362,-1.6685) -- (0.5780,-1.6994) -- (0.6179,-1.7325) -- (0.6560,-1.7677) -- (0.6920,-1.8051) -- (0.7260,-1.8443) -- (0.7577,-1.8854) -- (0.7871,-1.9282) -- (0.8140,-1.9725) -- (0.8385,-2.0182) -- (0.8604,-2.0653) -- (0.8797,-2.1135) -- (0.8963,-2.1626) -- (0.9101,-2.2126) -- (0.9212,-2.2633) -- (0.9294,-2.3146) -- (0.9348,-2.3662) -- (0.9373,-2.4180) -- (0.9369,-2.4699) -- (0.9337,-2.5217) -- (0.9276,-2.5732) -- (0.9187,-2.6243) -- (0.9069,-2.6749) -- (0.8924,-2.7247) -- (0.8752,-2.7736) -- (0.8552,-2.8215) -- (0.8327,-2.8683) -- (0.8076,-2.9137) -- (0.7800,-2.9576) -- (0.7500,-3.0000);\draw[object] (-0.75,-3)--(-.2,-3) (.2,-3)--(.75,-3);\draw[traj] (2.2500,-3.0000) -- (2.2200,-2.9576) -- (2.1924,-2.9137) -- (2.1673,-2.8683) -- (2.1448,-2.8215) -- (2.1248,-2.7736) -- (2.1076,-2.7247) -- (2.0931,-2.6749) -- (2.0813,-2.6243) -- (2.0724,-2.5732) -- (2.0663,-2.5217) -- (2.0631,-2.4699) -- (2.0627,-2.4180) -- (2.0652,-2.3662) -- (2.0706,-2.3146) -- (2.0788,-2.2633) -- (2.0899,-2.2126) -- (2.1037,-2.1626) -- (2.1203,-2.1135) -- (2.1396,-2.0653) -- (2.1615,-2.0182) -- (2.1860,-1.9725) -- (2.2129,-1.9282) -- (2.2423,-1.8854) -- (2.2740,-1.8443) -- (2.3080,-1.8051) -- (2.3440,-1.7677) -- (2.3821,-1.7325) -- (2.4220,-1.6994) -- (2.4638,-1.6685) -- (2.5071,-1.6400) -- (2.5520,-1.6140) -- (2.5983,-1.5904) -- (2.6457,-1.5695) -- (2.6943,-1.5512) -- (2.7438,-1.5357) -- (2.7941,-1.5229) -- (2.8450,-1.5129) -- (2.8964,-1.5057) -- (2.9481,-1.5014) -- (3.0000,-1.5000) -- (3.0519,-1.5014) -- (3.1036,-1.5057) -- (3.1550,-1.5129) -- (3.2059,-1.5229) -- (3.2562,-1.5357) -- (3.3057,-1.5512) -- (3.3543,-1.5695) -- (3.4017,-1.5904) -- (3.4480,-1.6140) -- (3.4929,-1.6400) -- (3.5362,-1.6685) -- (3.5780,-1.6994) -- (3.6179,-1.7325) -- (3.6560,-1.7677) -- (3.6920,-1.8051) -- (3.7260,-1.8443) -- (3.7577,-1.8854) -- (3.7871,-1.9282) -- (3.8140,-1.9725) -- (3.8385,-2.0182) -- (3.8604,-2.0653) -- (3.8797,-2.1135) -- (3.8963,-2.1626) -- (3.9101,-2.2126) -- (3.9212,-2.2633) -- (3.9294,-2.3146) -- (3.9348,-2.3662) -- (3.9373,-2.4180) -- (3.9369,-2.4699) -- (3.9337,-2.5217) -- (3.9276,-2.5732) -- (3.9187,-2.6243) -- (3.9069,-2.6749) -- (3.8924,-2.7247) -- (3.8752,-2.7736) -- (3.8552,-2.8215) -- (3.8327,-2.8683) -- (3.8076,-2.9137) -- (3.7800,-2.9576) -- (3.7500,-3.0000);\draw[object] (2.25,-3)--(2.8,-3) (3.2,-3)--(3.75,-3);
\draw[object] (-.2,0)--(-.2,-.62)--(3.2,-.62)--(3.2,0) (.2,0)--(.2,-.30)--(2.8,-.30)--(2.8,0);
\draw[object] (-.2,-3)--(-.2,-3.62)--(3.2,-3.62)--(3.2,-3) (.2,-3)--(.2,-3.30)--(2.8,-3.30)--(2.8,-3);
\node at (1.5,1.2){小于半球};\node at (1.5,-.9){两泡连通};\node at (1.5,-3.9){大于半球};\end{tikzpicture}\figcaption{图 热4.8　同一管口的球冠越高，曲率半径未必越大。}\end{minipage}\end{center}

\workarea

\clearpage

\pdfbookmark[2]{第 9 题 润湿浮球温度变化时浸入体积不变条件}{q9}

\section*{第 9 题\quad 润湿浮球温度变化时浸入体积不变条件}

半径$R$、密度$\rho$的均匀球浮在0℃液体表面，球心高出远处液面$h$。球的热膨胀忽略；液体体膨胀系数为$\alpha$，表面张力$\sigma(t)=\sigma_0(1-\beta t)$，$t$为摄氏温度。液体完全润湿球。采用图中的接触圆与液面几何，证明0℃附近温度小幅变化时，浸入体积一阶不变的条件为\[\frac{\beta-\alpha}{\alpha}=\frac{2R^4\rho g}{3\sigma_0(R^2-h^2)}.\]

\begin{center}\begin{minipage}{0.96\linewidth}\centering\begin{tikzpicture}[x=1cm,y=1cm,font=\small,every node/.append style={align=center}]
\pgfmathsetmacro{\rc}{sqrt(1.3*1.3-.6*.6)}\fill[blue!12] (-2,-1.8) rectangle (2,0);\draw[object] (-2,0)--(2,0);\draw[object] (0,.6) circle (1.3);
\draw[guide] (-\rc,0)--(\rc,0);\draw[guide] (0,.6)--(\rc,0);\node at (.85,.65){$R$};
\draw[<->] (-1.6,0)--(-1.6,.6) node[midway,left]{$h$};\draw[guide] (\rc,0)--(2.05,-.1);\node[right] at (2.05,-.1){接触圆半径$r_c$};
\draw[force] (0,.6)--(0,-1.15) node[right]{$mg$};\draw[force] (-2.1,-.9)--(-2.1,.5) node[left]{$\rho_lgV_L$};\end{tikzpicture}\figcaption{图 热4.9　图中浸入体积按远处液面与球的几何截切计算。}\end{minipage}\end{center}

\workarea

\clearpage

\pdfbookmark[2]{第 10 题 对称连体液膜中环形细线张力}{q10}

\section*{第 10 题\quad 对称连体液膜中环形细线张力}

两个相同球形液膜接触后形成对称连体膜：外侧是两个超过半球的球面，曲率半径均为$R$；中间为半径$r$的圆膜，边缘围一根均匀柔软细线。每个液膜都有两个表面，表面张力系数为$\sigma$，重力忽略。求细线内部张力$T$，并讨论柔软线不能受压所造成的几何限制。

\begin{center}\begin{minipage}{0.96\linewidth}\centering\begin{tikzpicture}[x=1cm,y=1cm,font=\small,every node/.append style={align=center}]\draw[traj] (0.0000,1.2000) -- (-0.0540,1.2387) -- (-0.1097,1.2749) -- (-0.1669,1.3086) -- (-0.2255,1.3398) -- (-0.2855,1.3684) -- (-0.3467,1.3942) -- (-0.4090,1.4174) -- (-0.4722,1.4377) -- (-0.5363,1.4552) -- (-0.6011,1.4699) -- (-0.6664,1.4817) -- (-0.7322,1.4906) -- (-0.7984,1.4966) -- (-0.8648,1.4996) -- (-0.9312,1.4997) -- (-0.9975,1.4968) -- (-1.0637,1.4910) -- (-1.1296,1.4823) -- (-1.1950,1.4707) -- (-1.2598,1.4562) -- (-1.3239,1.4389) -- (-1.3872,1.4187) -- (-1.4495,1.3957) -- (-1.5108,1.3700) -- (-1.5708,1.3416) -- (-1.6296,1.3106) -- (-1.6869,1.2770) -- (-1.7426,1.2410) -- (-1.7967,1.2024) -- (-1.8491,1.1616) -- (-1.8996,1.1184) -- (-1.9481,1.0730) -- (-1.9946,1.0256) -- (-2.0389,0.9761) -- (-2.0810,0.9247) -- (-2.1208,0.8716) -- (-2.1582,0.8167) -- (-2.1931,0.7602) -- (-2.2255,0.7022) -- (-2.2553,0.6428) -- (-2.2824,0.5822) -- (-2.3068,0.5204) -- (-2.3285,0.4576) -- (-2.3474,0.3939) -- (-2.3634,0.3294) -- (-2.3765,0.2643) -- (-2.3868,0.1987) -- (-2.3941,0.1327) -- (-2.3985,0.0664) -- (-2.4000,0.0000) -- (-2.3985,-0.0664) -- (-2.3941,-0.1327) -- (-2.3868,-0.1987) -- (-2.3765,-0.2643) -- (-2.3634,-0.3294) -- (-2.3474,-0.3939) -- (-2.3285,-0.4576) -- (-2.3068,-0.5204) -- (-2.2824,-0.5822) -- (-2.2553,-0.6428) -- (-2.2255,-0.7022) -- (-2.1931,-0.7602) -- (-2.1582,-0.8167) -- (-2.1208,-0.8716) -- (-2.0810,-0.9247) -- (-2.0389,-0.9761) -- (-1.9946,-1.0256) -- (-1.9481,-1.0730) -- (-1.8996,-1.1184) -- (-1.8491,-1.1616) -- (-1.7967,-1.2024) -- (-1.7426,-1.2410) -- (-1.6869,-1.2770) -- (-1.6296,-1.3106) -- (-1.5708,-1.3416) -- (-1.5108,-1.3700) -- (-1.4495,-1.3957) -- (-1.3872,-1.4187) -- (-1.3239,-1.4389) -- (-1.2598,-1.4562) -- (-1.1950,-1.4707) -- (-1.1296,-1.4823) -- (-1.0637,-1.4910) -- (-0.9975,-1.4968) -- (-0.9312,-1.4997) -- (-0.8648,-1.4996) -- (-0.7984,-1.4966) -- (-0.7322,-1.4906) -- (-0.6664,-1.4817) -- (-0.6011,-1.4699) -- (-0.5363,-1.4552) -- (-0.4722,-1.4377) -- (-0.4090,-1.4174) -- (-0.3467,-1.3942) -- (-0.2855,-1.3684) -- (-0.2255,-1.3398) -- (-0.1669,-1.3086) -- (-0.1097,-1.2749) -- (-0.0540,-1.2387) -- (0.0000,-1.2000);
\draw[traj] (-0.0000,1.2000) -- (0.0540,1.2387) -- (0.1097,1.2749) -- (0.1669,1.3086) -- (0.2255,1.3398) -- (0.2855,1.3684) -- (0.3467,1.3942) -- (0.4090,1.4174) -- (0.4722,1.4377) -- (0.5363,1.4552) -- (0.6011,1.4699) -- (0.6664,1.4817) -- (0.7322,1.4906) -- (0.7984,1.4966) -- (0.8648,1.4996) -- (0.9312,1.4997) -- (0.9975,1.4968) -- (1.0637,1.4910) -- (1.1296,1.4823) -- (1.1950,1.4707) -- (1.2598,1.4562) -- (1.3239,1.4389) -- (1.3872,1.4187) -- (1.4495,1.3957) -- (1.5108,1.3700) -- (1.5708,1.3416) -- (1.6296,1.3106) -- (1.6869,1.2770) -- (1.7426,1.2410) -- (1.7967,1.2024) -- (1.8491,1.1616) -- (1.8996,1.1184) -- (1.9481,1.0730) -- (1.9946,1.0256) -- (2.0389,0.9761) -- (2.0810,0.9247) -- (2.1208,0.8716) -- (2.1582,0.8167) -- (2.1931,0.7602) -- (2.2255,0.7022) -- (2.2553,0.6428) -- (2.2824,0.5822) -- (2.3068,0.5204) -- (2.3285,0.4576) -- (2.3474,0.3939) -- (2.3634,0.3294) -- (2.3765,0.2643) -- (2.3868,0.1987) -- (2.3941,0.1327) -- (2.3985,0.0664) -- (2.4000,0.0000) -- (2.3985,-0.0664) -- (2.3941,-0.1327) -- (2.3868,-0.1987) -- (2.3765,-0.2643) -- (2.3634,-0.3294) -- (2.3474,-0.3939) -- (2.3285,-0.4576) -- (2.3068,-0.5204) -- (2.2824,-0.5822) -- (2.2553,-0.6428) -- (2.2255,-0.7022) -- (2.1931,-0.7602) -- (2.1582,-0.8167) -- (2.1208,-0.8716) -- (2.0810,-0.9247) -- (2.0389,-0.9761) -- (1.9946,-1.0256) -- (1.9481,-1.0730) -- (1.8996,-1.1184) -- (1.8491,-1.1616) -- (1.7967,-1.2024) -- (1.7426,-1.2410) -- (1.6869,-1.2770) -- (1.6296,-1.3106) -- (1.5708,-1.3416) -- (1.5108,-1.3700) -- (1.4495,-1.3957) -- (1.3872,-1.4187) -- (1.3239,-1.4389) -- (1.2598,-1.4562) -- (1.1950,-1.4707) -- (1.1296,-1.4823) -- (1.0637,-1.4910) -- (0.9975,-1.4968) -- (0.9312,-1.4997) -- (0.8648,-1.4996) -- (0.7984,-1.4966) -- (0.7322,-1.4906) -- (0.6664,-1.4817) -- (0.6011,-1.4699) -- (0.5363,-1.4552) -- (0.4722,-1.4377) -- (0.4090,-1.4174) -- (0.3467,-1.3942) -- (0.2855,-1.3684) -- (0.2255,-1.3398) -- (0.1669,-1.3086) -- (0.1097,-1.2749) -- (0.0540,-1.2387) -- (-0.0000,-1.2000);
\draw[object,very thick] (0,-1.2)--(0,1.2);\node[right] at (.1,0){圆膜截面};\draw[guide] (-.9,0)--(0,1.2) node[midway,left]{$R$};
\node[above] at (0,1.6){中间圆膜半径$r$};\end{tikzpicture}\figcaption{图 热4.10　截面中两外球冠相同，中间膜为圆形平面。}\end{minipage}\end{center}

\workarea

\clearpage

\pdfbookmark[2]{第 11 题 润湿竖直平板上的毛细上升高度}{q11}

\section*{第 11 题\quad 润湿竖直平板上的毛细上升高度}

一块竖直无限大平板部分浸入密度$\rho$、表面张力系数$\sigma$的液体，与液体的接触角为$\theta$，$0\le\theta\le\pi/2$。远处液面水平，外界气压均匀。求液体沿板上升的高度$h$。要求从液面曲率与静水压平衡推导，不把平板当作圆毛细管。

\begin{center}\begin{minipage}{0.96\linewidth}\centering\begin{tikzpicture}[x=1cm,y=1cm,font=\small,every node/.append style={align=center}]\fill[blue!12] (0,-.7)--(6,-.7)--(6,0)--(0,0)--cycle;\draw[object,very thick] (0,-.7)--(0,3);
\draw[traj] (0,2.2) .. controls (.2,.6) and (1.8,.02) .. (5.8,0);\draw[guide] (0,0)--(6,0);
\draw[<->] (-.4,0)--(-.4,2.2) node[midway,left]{$h$};\node at (1,2.7){接触角$\theta$};\node at (3,-.4){远处水平液面$y=0$};\end{tikzpicture}\figcaption{图 热4.11　沿板方向平移不改变液面，第二主曲率为零。}\end{minipage}\end{center}

\workarea

\clearpage

\pdfbookmark[2]{第 12 题 一维晶体作用势、弹性和断裂形变}{q12}

\section*{第 12 题\quad 一维晶体作用势、弹性和断裂形变}

均匀一维晶体由相同原子组成，只考虑最近邻作用，势能$V(x)=-A/x^6+B/x^{12}$，$A,B>0$，$x$为相邻间距。求零温静态模型下的平衡距离$x_0$；求小相对形变$\varepsilon=(x-x_0)/x_0$对应的力系数$K$，即外拉力$F\simeq K\varepsilon$；在缓慢增加外拉力的模型下，求失去稳定均匀平衡时的临界相对形变。分别注明相对形变系数、键刚度和实际有限形变的区别。

\begin{center}\begin{minipage}{0.96\linewidth}\centering\begin{tikzpicture}[x=1cm,y=1cm,font=\small,every node/.append style={align=center}]\foreach \x in {0,1.5,3,4.5,6}{\filldraw[fill=blue!15] (\x,0) circle(.18);}
\foreach \x in {0,1.5,3,4.5}{\draw[object] (\x+.18,0)--(\x+1.32,0);}
\draw[<->] (1.5,.65)--(3,.65) node[midway,above]{$x$};\draw[force] (-.2,0)--(-1,0) node[left]{$F$};
\draw[force] (6.2,0)--(7,0) node[right]{$F$};\node at (3,-.75){只计最近邻键；相对形变$\varepsilon=(x-x_0)/x_0$};\end{tikzpicture}\figcaption{图 热4.12　每条串联键承受同一外拉力。}\end{minipage}\end{center}

\workarea

\clearpage

\pdfbookmark[2]{第 13 题 NaCl晶体结合能的估算}{q13}

\section*{第 13 题\quad NaCl晶体结合能的估算}

NaCl晶体平衡时最近邻异号离子距离$r_0=2.81\times10^{-10}\,\mathrm m$，绝热压缩系数$k_s=3.3\times10^{-11}\,\mathrm{Pa^{-1}}$，马德隆常量$\alpha_M=1.75$。采用每个离子对势能$u(r)=B/r^n-\alpha_M e^2/(4\pi\varepsilon_0r)$的库仑吸引与幂次排斥模型，晶体体积$V=2Nr^3$，热运动能量对弹性和结合能的影响忽略。求1 mol NaCl的平衡势能$E_{p0}$及把晶体拆成孤立离子所需的结合能大小。

\begin{center}\begin{minipage}{0.96\linewidth}\centering\begin{tikzpicture}[x=1cm,y=1cm,font=\small,every node/.append style={align=center}]\foreach \i in {0,1,2,3}{\foreach \j in {0,1,2,3}{\pgfmathtruncatemacro{\parity}{mod(\i+\j,2)}
\ifnum\parity=0 \node[draw,circle,fill=blue!12,inner sep=3pt] at(\i,\j){$+$};\else\node[draw,circle,fill=brown!12,inner sep=3pt] at(\i,\j){$-$};\fi}}
\draw[<->] (0,-.6)--(1,-.6) node[midway,below]{$r_0$};\node[right] at (3.7,1.5){交替离子截面\\库仑项与排斥项};\end{tikzpicture}\figcaption{图 热4.13　马德隆常量计入整个晶格的静电作用。}\end{minipage}\end{center}

\workarea

\clearpage

\pdfbookmark[2]{第 14 题 等温压缩氟利昂的液化质量}{q14}

\section*{第 14 题\quad 等温压缩氟利昂的液化质量}

质量$M=2.0\,\mathrm{kg}$的氟利昂在$-13{}^\circ\mathrm{C}$下由$V_i=0.19\,\mathrm{m^3}$等温压缩至$V_f=0.10\,\mathrm{m^3}$。摩尔质量$\mu=121\times10^{-3}\,\mathrm{kg/mol}$，该温度液体密度$\rho_l=1.44\times10^3\,\mathrm{kg/m^3}$，饱和蒸汽压$p_s=2.08\times10^5\,\mathrm{Pa}$。饱和蒸汽视为理想气体。求压缩中液化的质量；先检查初态是否已有液体，保留终态液体所占体积。

\begin{center}\begin{minipage}{0.96\linewidth}\centering\begin{tikzpicture}[x=1cm,y=1cm,font=\small,every node/.append style={align=center}]\draw[object] (0,0) rectangle(2.3,2.8);\draw[object,very thick] (0,2.5)--(2.3,2.5);\node at(1.15,1.3){初态\\$V_i=0.19$};
\draw[motion] (2.7,1.4)--(3.7,1.4);\draw[object] (4.2,0) rectangle(6.5,2.8);\draw[object,very thick] (4.2,1.9)--(6.5,1.9);
\fill[blue!18] (4.2,0) rectangle(6.5,.5);\node at(5.35,1.2){饱和蒸汽};\node at(5.35,.25){液体};\node[below] at(3.2,-.4){全程$T=260.15\,\mathrm K$};\end{tikzpicture}\figcaption{图 热4.14　气、液两相共同占据终态体积。}\end{minipage}\end{center}

\workarea

\clearpage

\pdfbookmark[2]{第 15 题 互不相溶液体共沸温度与蒸汽组成}{q15}

\section*{第 15 题\quad 互不相溶液体共沸温度与蒸汽组成}

互不相溶液体A、B的饱和蒸汽压满足$\ln(p_i/p_0)=a_i/T+b_i$，$T$用绝对温标，$p_0$为标准大气压。测得数据：
\begin{center}\begin{tabular}{ccc}\toprule 温度 & $p_A/p_0$ & $p_B/p_0$\\\midrule40℃&0.284&0.07278\\90℃&1.476&0.6918\\\bottomrule\end{tabular}\end{center}
外压为$p_0$，先求A、B各自沸点。再将100 g A、100 g B置于图示容器，在B上方覆盖不挥发、与A和B均不相溶的薄液层C，抑制B在自由液面的蒸发；各层静水压忽略。A、B蒸汽在界面气泡中混合，服从道尔顿分压定律，摩尔质量比$\mu_A/\mu_B=8$。持续缓慢加热，温度图有两段平台。求第一、第二平台温度$t_1,t_2$（到1℃），及图中第一平台结束、升温至第二平台的阶段内$\tau_1$时刻两液体剩余质量（到0.1 g）。气相滞留量及非沸腾阶段的自由蒸发损失忽略。

\begin{center}\begin{minipage}{0.96\linewidth}\centering\begin{tikzpicture}[x=1cm,y=1cm,font=\small,every node/.append style={align=center}]\draw[object] (0,0) rectangle(2.4,3);\fill[blue!18] (0,0) rectangle(2.4,1);\fill[brown!18] (0,1) rectangle(2.4,2);
\fill[gray!18] (0,2) rectangle(2.4,2.3);\node at(1.2,.5){A};\node at(1.2,1.5){B};\node at(1.2,2.15){C};
\draw[motion] (1.2,1)--(1.2,2.8) node[above]{混合气泡};
\draw[axis] (3.5,0)--(8.5,0) node[right]{$\tau$};\draw[axis] (3.5,0)--(3.5,3.3) node[above]{$t$};
\draw[traj] (3.6,.2)--(4.6,1.1)--(6.1,1.1)--(7.3,2.6)--(8.2,2.6);
\node[left] at(3.5,1.1){$t_1$};\node[left] at(3.5,2.6){$t_2$};\draw[guide] (6.65,0) node[below]{$\tau_1$}--(6.65,1.8);\end{tikzpicture}\figcaption{图 热4.15　A、B共同沸腾形成第一平台，剩余液体形成第二平台。}\end{minipage}\end{center}

\workarea

\clearpage

\pdfbookmark[2]{第 16 题 冰水水蒸气三相定容加热后的质量分配}{q16}

\section*{第 16 题\quad 冰水水蒸气三相定容加热后的质量分配}

密闭刚性容器内只有冰、水、水蒸汽，初态三相质量均为1 g，在三相点$t=0.01{}^\circ\mathrm{C}$、$p=4.58\,\mathrm{mmHg}$平衡。保持容积不变，缓慢输入热量$Q=0.255\,\mathrm{kJ}$。求终态三相质量的估计。升华热$L_s=2.83\,\mathrm{kJ/g}$，汽化热$L_v=2.49\,\mathrm{kJ/g}$；水蒸汽视为理想气体，取水、冰密度约$10^3,0.9\times10^3\,\mathrm{kg/m^3}$。先检查三相是否仍共存，再说明忽略凝聚相体积变化的误差。

\begin{center}\begin{minipage}{0.96\linewidth}\centering\begin{tikzpicture}[x=1cm,y=1cm,font=\small,every node/.append style={align=center}]\draw[object] (0,0) rectangle(5.5,3.2);\fill[blue!10] (0,0) rectangle(5.5,.5);
\filldraw[fill=gray!25] (.3,.15) rectangle(1.1,.85);\node at(3,2){水蒸汽1 g\\占绝大部分容积};
\node[below] at(.7,-.1){冰1 g};\node[below] at(3.8,-.1){水1 g};\draw[force] (6.5,1)--(5.5,1) node[midway,above]{$Q$};\end{tikzpicture}\figcaption{图 热4.16　容器刚性密闭，三相共存时温度和压强固定。}\end{minipage}\end{center}

\workarea

\clearpage

\pdfbookmark[2]{第 17 题 地幔硅酸盐熔点随深度变化}{q17}

\section*{第 17 题\quad 地幔硅酸盐熔点随深度变化}

地幔固态硅酸盐的熔化潜热为单位质量的$L_m$，液、固密度比$\alpha=\rho_l/\rho_s$。$r$表示离地心的距离。结合静力平衡和克拉珀龙方程，求熔点$T_m$随$r$的变化率。先以一般重力加速度$g_r$表示，再在均匀密度地球模型中写成表面重力$g$、地球半径$R$的形式。上部地幔某处$T_m=1300{}^\circ\mathrm{C}$、$\alpha=0.9$、$L_m=10^5\,\mathrm{cal/kg}$，估算向地心深入1000 m后的熔点变化，取$g=9.80\,\mathrm{m/s^2}$、$1\,\mathrm{cal}=4.184\,\mathrm J$。

\begin{center}\begin{minipage}{0.96\linewidth}\centering\begin{tikzpicture}[x=1cm,y=1cm,font=\small]
\draw[object] (0,0) circle(1.45);
\fill[blue!18] (10:1.1) arc[start angle=10,end angle=30,radius=1.1]--(30:1.2) arc[start angle=30,end angle=10,radius=1.2]--cycle;
\draw[guide] (0,0)--(20:1.15) node[midway,above] {$r$};
\node[below left] at (0,0) {地心};
\draw[object,fill=blue!10] (4,-.7) rectangle(5.2,.7);
\draw[force] (2.5,0)--(4,0) node[midway,above] {$p(r)S$};
\draw[force] (6.5,0)--(5.2,0) node[midway,above] {$p(r+\dd r)S$};
\draw[force] (4.95,-.45)--(4.15,-.45) node[midway,below=4pt] {$\rho_sS\dd r\,g_r$};
\draw[<->] (4,1.1)--(5.2,1.1) node[midway,above] {$\dd r$};
\draw[axis] (3,-1.4)--(6.4,-1.4) node[right] {$r$向外};
\node at (4.6,1.8) {局部薄层受力};
\end{tikzpicture}\figcaption{图 热4.17　径向坐标向外增加，重力指向地心。}\end{minipage}\end{center}

\workarea

\end{document}
